% Read STYLE.md before editing prose in this file.
% Follow the Content Creation Workflow in WORKFLOW.md §Content Creation Workflow.

\section{Introduction}
\label{sec:intro}

Dolphins have long been recognized as among the most communicatively
sophisticated non-human animals on Earth~\citep{Wonko2021communication}.
Their vocalizations span a wide frequency range---from low-frequency
burst-pulse sounds near 200\Hz{} to high-frequency whistles exceeding
20\kHz{}---and serve functions including individual identification,
group coordination, and echolocation~\citep{Trillian2023signal}.
Yet despite this richness, one category of dolphin vocalization has
attracted remarkably little systematic study: the acoustic signals
produced immediately before a group or individual permanently departs a
location.

\begin{hypothesis}[Farewell Specificity]
  \label{hyp:specificity}
  Pre-departure vocalizations in \spname{Tursiops truncatus} form a
  structurally distinct class, separable from non-departure vocalizations
  by measurable acoustic features including frequency contour,
  inter-pulse interval, and burst duration.
\end{hypothesis}

This paper tests \cref{hyp:specificity} by applying computational
linguistics methods---specifically hidden Markov model (HMM) sequence
alignment and spectral envelope analysis---to a curated corpus of
42~departure events recorded across three Atlantic coastal populations
between 2021 and 2024~\citep{GuideResearch2025dataset}.

\subsection{Contributions}
\label{subsec:contrib}

The main contributions of this work are as follows.

\begin{itemize}
  \item We introduce the \emph{Atlantic Cetacean Farewell Corpus}
        (ACFC v1.0), the first annotated dataset of departure-specific
        \spname{T.\ truncatus} vocalizations.

  \item We identify three recurring sequence archetypes---which we term
        \aseq{brief-farewell}, \aseq{extended-farewell}, and
        \aseq{urgent-departure}---and characterize their acoustic profiles
        (see \cref{tab:vocalization-comparison}).

  \item We demonstrate that a bigram HMM trained on ACFC achieves
        86.4\% accuracy in distinguishing departure from non-departure
        sequences, substantially above the chance baseline of 50\%.
\end{itemize}

\subsection{Paper Organization}
\label{subsec:organization}

\Cref{sec:background} reviews related literature on cetacean
communication and computational approaches to bioacoustic classification.
\Cref{sec:methodology} describes the data collection protocol,
annotation scheme, and computational methods.
We conclude with a discussion of implications for interspecies
communication theory.
